\documentclass[11pt,a4paper]{article}
\usepackage[utf8]{inputenc}
\usepackage[margin=1in]{geometry}
\usepackage{amsmath,amssymb,amsthm}
\usepackage{listings}
\usepackage{xcolor}
\usepackage{hyperref}

\newtheorem{theorem}{Theorem}
\newtheorem{lemma}[theorem]{Lemma}

\lstset{
  language=C++,
  basicstyle=\ttfamily\small,
  keywordstyle=\color{blue}\bfseries,
  commentstyle=\color{green!60!black},
  stringstyle=\color{red!70!black},
  numbers=left,
  numberstyle=\tiny\color{gray},
  breaklines=true,
  frame=single,
  tabsize=4,
  showstringspaces=false
}

\title{IOI 2011 -- Crocodile}
\author{Solution}
\date{}

\begin{document}
\maketitle

\section{Problem Summary}
A network of $N$ chambers connected by $M$ bidirectional corridors with travel times. Some $K$ chambers are exits. A person starts at chamber~$0$. After choosing which corridor to take, a crocodile may block it, forcing the person to take the second-best option. Find the minimum guaranteed escape time from chamber~$0$ under optimal adversarial play by the crocodile.

\section{Solution}

\subsection{Key Insight}
Since the crocodile always blocks the best option, the guaranteed escape cost from any node equals the \emph{second-shortest} distance to an exit. We track the two shortest distances to any exit for each node.

\subsection{Algorithm: Modified Multi-Source Dijkstra}
\begin{enumerate}
    \item Initialize all exits with $d_1 = d_2 = 0$.
    \item Use a min-heap keyed on $d_2$. A node is ``finalized'' when its $d_2$ is determined.
    \item When finalizing node~$u$ (with second-best distance $d_2[u]$), for each neighbor~$v$ via edge of weight~$w$: the candidate distance is $d_2[u] + w$. Update $d_1[v]$ and $d_2[v]$ accordingly:
    \begin{itemize}
        \item If $d_2[u] + w < d_1[v]$: shift $d_1[v]$ to $d_2[v]$ and set $d_1[v] = d_2[u] + w$.
        \item Else if $d_2[u] + w < d_2[v]$: set $d_2[v] = d_2[u] + w$.
    \end{itemize}
    \item When $d_2[v]$ improves, push $(d_2[v], v)$ into the heap.
\end{enumerate}

The answer is $d_2[0]$.

\subsection{Correctness}

\begin{theorem}
Using $d_2[u]$ (the second-best distance) for relaxation correctly computes the guaranteed escape cost.
\end{theorem}
\begin{proof}
From node~$u$, the person's optimal strategy is to always head toward the best available corridor leading to an exit. The crocodile blocks the best choice, forcing the person to use the second-best option. Hence the guaranteed cost from $u$ is $d_2[u]$.

When relaxing from $u$ to $v$, arriving at $u$ costs at least $d_2[u]$ (guaranteed), plus $w$ to traverse the edge. This candidate replaces $d_1[v]$ or $d_2[v]$ if it improves either. Dijkstra's monotonicity ensures that when $d_2[u]$ is finalized, no shorter path to $u$ via two distinct routes exists.
\end{proof}

\subsection{Complexity}
\begin{itemize}
    \item \textbf{Time:} $O((N + M) \log N)$.
    \item \textbf{Space:} $O(N + M)$.
\end{itemize}

\section{Implementation}

\begin{lstlisting}
#include <bits/stdc++.h>
using namespace std;

int travel_plan(int N, int M, int R[][2], int L[], int K, int P[]) {
    vector<vector<pair<int,int>>> adj(N);
    for (int i = 0; i < M; i++) {
        adj[R[i][0]].push_back({R[i][1], L[i]});
        adj[R[i][1]].push_back({R[i][0], L[i]});
    }

    const long long INF = 1e18;
    vector<array<long long, 2>> d(N, {INF, INF});

    priority_queue<pair<long long,int>, vector<pair<long long,int>>,
                   greater<>> pq;

    for (int i = 0; i < K; i++) {
        d[P[i]] = {0, 0};
        pq.push({0, P[i]});
    }

    vector<bool> done(N, false);

    while (!pq.empty()) {
        auto [dist, u] = pq.top(); pq.pop();
        if (done[u]) continue;
        if (dist > d[u][1]) continue;
        done[u] = true;

        for (auto [v, w] : adj[u]) {
            long long nd = d[u][1] + w;
            if (nd < d[v][0]) {
                d[v][1] = d[v][0];
                d[v][0] = nd;
            } else if (nd < d[v][1]) {
                d[v][1] = nd;
            } else {
                continue;
            }
            if (d[v][1] < INF)
                pq.push({d[v][1], v});
        }
    }

    return (int)d[0][1];
}
\end{lstlisting}

\end{document}
